\section{Appendix: proofs and verification scope}

\subsection{Fixed-law bounds and compatibility}

The fixed released-law calculation precedes the optimization over released laws. We attribute its mathematical antecedent to \citet[Theorem 3.2 and eq. (3.1)]{FanShermanShum2014}. The formal declaration and proof displayed here are independently checked under the paper's stated model conditions. The compatibility criterion in \cref{obj:prop:compatibility} identifies the released laws to which the calculation applies. The first auxiliary result records the arm-label masses implied by a score law and release rule.

\begin{lemmav}{lem:released-arm-cell-masses}[Released arm-cell masses]
Let \(H\) be a measure on the score interval \(\mathcal E\), let \(g:\mathcal E\to\mathcal R_J\) be a release rule, and let \(P\in\mathfrak P(H,g)\) as in \cref{obj:def:causal-law-class}. Write \(P_{\mathrm{rel}}=\mathcal L_P(R,A,Y)\) for the released law. Then \(H\) and \(P_{\mathrm{rel}}\) are probability laws, and, for every \(a\in\mathcal A\) and \(r\in\mathcal R_J\),
\[
P_{\mathrm{rel}}(A=a,R=r)=q_{ar},
\]
where \(q_{ar}\) is the arm-cell mass defined in \cref{obj:def:arm-cell-primitives}.
\end{lemmav}

These mass equalities connect the causal-law class to the compatibility criterion in \cref{obj:prop:compatibility}. For each compatible released law, the fixed-law result gives sharp bounds through quantile rearrangement within positive-mass arm-label cells.

\begin{lemmav}{lem:fixed-released-law-baseline}[Sharp fixed released law bounds]
Let \(H\) be a measure on the score interval \(\mathcal E=[\varepsilon,1-\varepsilon]\), let \(g:\mathcal E\to\mathcal R_J\) be a release rule, and let \(P_{\mathrm{rel}}\) be a measure on the released label, treatment arm, and outcome. Suppose:

\begin{itemize}
\item \textbf{(Overlap.)} The score support satisfies \cref{obj:ass:overlap}.
\item \textbf{(Release.)} The rule \(g\) is measurable.
\item \textbf{(Compatibility.)} \(\mathfrak M(H,g,P_{\mathrm{rel}})\ne\varnothing\).
\end{itemize}

For each arm \(a\in\{0,1\}\) and label \(r\), let \(C_r=\{e:g(e)=r\}\), let \(p_a(e)\) be the probability of arm \(a\) at score \(e\), and set \(q_{ar}=\int_{C_r}p_a(e)\,H(de)\). In each cell with \(q_{ar}>0\), let \(F_{ar}\) be the outcome law conditional on arm \(a\) and label \(r\) under \(P_{\mathrm{rel}}\), let \(E_{ar}\) have law \(q_{ar}^{-1}\mathbf 1_{C_r}(e)p_a(e)\,H(de)\), and let \(W_{ar}=1/p_a(E_{ar})\). Writing \(Q\) for a generalized quantile, define
\[
\underline\mu_a
=\sum_{r:q_{ar}>0}q_{ar}\int_0^1 Q_{F_{ar}}(u)Q_{W_{ar}}(1-u)\,du,
\qquad
\overline\mu_a
=\sum_{r:q_{ar}>0}q_{ar}\int_0^1 Q_{F_{ar}}(u)Q_{W_{ar}}(u)\,du.
\]
Then, for every \(a\in\{0,1\}\),
\[
\left\{\mathbb E_P[Y_a]:P\in\mathfrak M(H,g,P_{\mathrm{rel}})\right\}
=[\underline\mu_a,\overline\mu_a].
\]
The attainable average treatment effects form the interval defined in \cref{obj:def:sharp-ate-set}:
\[
\left\{\mathbb E_P[Y_1-Y_0]:P\in\mathfrak M(H,g,P_{\mathrm{rel}})\right\}
=I(P_{\mathrm{rel}};H,g)
=[\underline\mu_1-\overline\mu_0,\,
  \overline\mu_1-\underline\mu_0].
\]
Moreover, there exist \(P^-,P^+\in\mathfrak M(H,g,P_{\mathrm{rel}})\) such that
\[
\mathbb E_{P^-}[Y_1-Y_0]=\underline\mu_1-\overline\mu_0,
\qquad
\mathbb E_{P^+}[Y_1-Y_0]=\overline\mu_1-\underline\mu_0.
\]
\end{lemmav}

The attainable endpoints make the interval sharp for each compatible released law. With the compatibility characterization in \cref{obj:prop:compatibility}, this fixed-law result supplies the interval whose width is optimized in \cref{obj:def:worst-case-ambiguity}.

\subsection{Ambiguity and finite-label design}

The all-label result in \cref{obj:thm:all-label-ambiguity} applies that outer optimization to the fixed-law interval, while \cref{obj:thm:universal-point-identification} characterizes universal point identification. For finite-label design, the proof first expresses the paired distortion as an optimization over arbitrary measurable partitions. The next lemma establishes the high-resolution limit over arbitrary measurable partitions and constructs an asymptotically attaining sequence of ordered interval cells.



\begin{lemmav}{lem:paired-high-resolution-quantization}[Paired high-resolution quantization]
Let \(\mathcal I=[a,b]\), let \(S\) be an integer, and let \(\beta_s:\mathcal I\times\mathcal I\to\mathbb R\), \(s=1,\ldots,S\), satisfy:
\begin{itemize}
\item \textbf{(Interval.)} \(a<b\).
\item \textbf{(Number of components.)} \(S\geq1\).
\item \textbf{(Continuity.)} Each \(\beta_s\) is continuous on \(\mathcal I\times\mathcal I\).
\item \textbf{(Positivity.)} \(\beta_s(x,z)>0\) for every \(s\) and every \(x,z\in\mathcal I\).
\end{itemize}
For each positive integer \(k\), define
\[
V_k=
\inf_{\substack{(B_1,\ldots,B_k)\ \mathrm{measurable,\ pairwise\ disjoint}\\
\bigcup_{j=1}^{k}B_j=\mathcal I,\quad z_{js}\in\mathcal I}}
\sum_{j=1}^{k}\sum_{s=1}^{S}
\int_{B_j}\beta_s(x,z_{js})|x-z_{js}|\,dx,
\qquad
w(x)=\sum_{s=1}^{S}\beta_s(x,x).
\]
Then
\[
\lim_{k\to\infty}kV_k
=\frac14\left(\int_a^b\sqrt{w(x)}\,dx\right)^2.
\]
For each positive integer \(k\), let \(a=t_0<\cdots<t_k=b\) be boundaries satisfying
\[
\int_a^{t_j}\sqrt{w(x)}\,dx
=\frac{j}{k}\int_a^b\sqrt{w(x)}\,dx
\qquad (j=0,\ldots,k).
\]
Set \(B_j=[t_{j-1},t_j)\) for \(j<k\), \(B_k=[t_{k-1},t_k]\), and \(z_{js}=(t_{j-1}+t_j)/2\) for every \(s\). These cells form a measurable, pairwise-disjoint partition of \(\mathcal I\), every \(z_{js}\) belongs to \(\mathcal I\), and
\[
\lim_{k\to\infty}
k\sum_{j=1}^{k}\sum_{s=1}^{S}
\int_{B_j}\beta_s(x,z_{js})|x-z_{js}|\,dx
=\frac14\left(\int_a^b\sqrt{w(x)}\,dx\right)^2.
\]
\end{lemmav}

The auxiliary result identifies the leading constant for the paired loss and gives an attaining sequence of ordered cells. It supplies the quantization component of the high-resolution conclusion in \cref{obj:thm:sharp-high-resolution-constant}.

\subsection{External-log excess risk}

The external-log analysis has two sampling contributions: released trial rows and independent score draws. The following certificate supplies lower bounds for both contributions to the excess-risk criterion in \cref{obj:def:external-excess-risk}.

\begin{lemmav}{lem:external-two-source-lower-certificate}[Two-source excess-risk lower bounds]
Let \(\mathcal E=[\varepsilon,1-\varepsilon]\) be the score interval, let \(\mathcal R_J\) be the finite label set, and fix a measurable release \(g:\mathcal E\to\mathcal R_J\). Let \(\mathbb Q_{P,H}^{n,m}\) be a family of joint sampling laws on \(n\) released trial rows and \(m\) score-log draws. Suppose:

\begin{itemize}
\item \textbf{(Overlap and level.)} \cref{obj:ass:overlap} holds and \(0<\alpha<1/2\).
\item \textbf{(Trial sampling.)} For every \(n,m\) and every admissible causal-law and score-law pair \((P,H)\) under \(g\), the trial marginal of \(\mathbb Q_{P,H}^{n,m}\) is the \(n\)-fold product of the released-row law under \(P\).
\item \textbf{(Score-log sampling.)} The joint sampling laws satisfy \cref{obj:ass:external-log-iid,obj:ass:external-log-independent} for every \(n,m\).
\end{itemize}

Set
\[
L_\alpha=1+(1-2\alpha)^2,\qquad
S_g=\sup_{\substack{r\in\mathcal R_J,\ e_1,e_2,e_3\in\mathcal E\\
e_1<e_2<e_3,\ g(e_1)=g(e_2)=g(e_3)=r}}
\frac{(e_3-e_2)(1-e_1/e_2)}
{\sqrt{(e_3-e_2)^2+(e_3-e_1)^2+(e_2-e_1)^2}},
\]
and
\[
c_{\mathrm{tr}}(\alpha)
=\frac{(1-2\alpha)\sqrt{\log L_\alpha}}{4},
\qquad
c_{\mathrm{log}}(g,\alpha)
=\frac{(1-2\alpha)\sqrt{\log L_\alpha}}{2\sqrt{3}}S_g.
\]
Then \(0<S_g\leq 1\). The external-score minimax excess risk of \cref{obj:def:external-excess-risk} satisfies, for every \(n,m\geq 1\),
\[
\mathcal R^\star_{n,m,\alpha}(g)
\geq \frac{c_{\mathrm{tr}}(\alpha)}{\sqrt n}.
\]
It also satisfies
\[
\liminf_{m\to\infty}\inf_{n\geq 1}
\left\{\sqrt m\,\mathcal R^\star_{n,m,\alpha}(g)\right\}
\geq c_{\mathrm{log}}(g,\alpha)>0.
\]
For every \(\rho>0\), the normalized risk of \cref{obj:def:external-normalized-risk} satisfies
\[
\liminf_{\substack{n,m\to\infty\\ n/m\to\rho}}
T_{n,m,\alpha}(g)
\geq
\max\left\{
\frac{c_{\mathrm{tr}}(\alpha)}{1+\sqrt\rho},
\frac{c_{\mathrm{log}}(g,\alpha)\sqrt\rho}{1+\sqrt\rho}
\right\}>0.
\]
\end{lemmav}

The trial bound applies at every pair of positive sample sizes. The score-log bound persists when trial size is unrestricted, and the final bound records both contributions at a fixed positive sample-size ratio. These bounds provide the lower-risk component of \cref{obj:thm:external-score-root-order}.

\subsection{Verification scope}

The displayed theorem and lemma declarations and their rendered proofs are machine-checked with Lean \texttt{4.33.0} at source commit \texttt{9135ac11b7d3e6ebd997f8c56bc48396dde396e4}. The machine-readable manifest \texttt{verification\_contract.json} maps each paper object identifier to its Lean declaration, source file, statement and proof audits, external dependencies, and sorry-free status. It records 44 displayed objects and 82 referenced Lean declarations. Objects marked \texttt{matched} have a checked formal counterpart; synthesized presentation definitions organize notation and are identified separately in that manifest. The fixed-law result is credited to \citet[Theorem 3.2 and eq. (3.1)]{FanShermanShum2014} as prior mathematics, while \cref{obj:lem:fixed-released-law-baseline} has no external formal dependency and its present proof is checked under the displayed assumptions.
